\section{SoCPilot Framework}
As illustrated in Figure~\ref{fig: overview}, SoCPilot translates natural-language design requirements and a C/C++ application into a customized SoC design. The framework contains four stages. First, it decomposes the requirements into a structured specification and selects reusable system IPs. Second, it profiles the application, evaluates candidate hardware implementations, and determines the hardware/software partition. Third, it generates and integrates the selected accelerators and refines the resulting SoC through system-level design space exploration (DSE) using latency and estimated area. Finally, it invokes the ESP implementation flow for FPGA synthesis, bitstream generation, and board-level validation. The LLM acts as the orchestration layer across these hardware and system-design stages, while profiling, HLS, simulation, and synthesis tools provide quantitative evidence for its decisions. The benchmark-specific embedded runtime code is manually implemented and treated as fixed infrastructure rather than an LLM-generated artifact.

